% Titre sans numéro de chapitre, à la même place
\renewcommand{\chaptername}{\phantom{Chapitre}}\renewcommand{\thechapter}{}
\chapter{Les structures de contrôles}
\renewcommand{\thesection}{\arabic{section}}
\section{Les boucles}
Les boucles sont une autre forme de structure de contrôle qui vont nous permettre de répéter des mêmes instructions.
On pourra répéter les instructions un nombre de fois fixe ou jusqu'à qu'une condition soit respectée.

\subsection{for}
La boucle for ("pour" en anglais) permet d'énumérer les éléments d'un ensemble (on parle d'\motcle{itérable}) et d'appliquer la même suite d'instruction à tous les éléments de cet ensemble.
C'est très utile, car ça permet de ne pas avoir à recopier le code pour chaque élément !

\begin{lstlisting}[style=pythonstyle, language=python]
for elem in iterable :
    cette ligne sera répéter pour chaque élément dans l itérable
    celle-ci aussi
celle-là ne sera executée qu une fois
\end{lstlisting}
\begin{center}\textbf{Comme pour les conditions, l'indentation nous permet de savoir ce qui est dans la boucle et ce qui ne l'est pas.}\end{center}

Notre première structure itérable sera l'intervalle des entiers relatif (positif et négatif), pour cela en Python, nous utiliserons la fonction~\motcle{range}. Voyons comment elle fonctionne, il y a trois manières d'utiliser la fonction range :
    \begin{center}
        Attention en informatique, on démarre toujours de 0.
    \end{center}
\begin{description}
    \item[range(fin)]: crée une liste de $fin$ nombres allant de $0$ à $fin-1$. Par exemple, range(4) permet de créer la suite 0, 1, 2, 3.
    \item[range(debut, fin)]: crée une liste de $fin - debut$ nombres allant de $debut$ à $fin-1$. Par exemple, range(4, 9) permet de créer la suite 4, 5, 6, 7, 8.
    \item[range(debut, fin, pas)]: crée une liste de $fin - debut$ nombres allant de $debut$ à au maximum $fin-1$ en comptant avec un$pas$. Par exemple, range(4, 9, 2) permet de créer la suite 4, 6, 8.\\
    D'ailleurs, il est possible de mettre un pas négatif, dans ce cas-là, on compte à l'envers. Par exemple, range(3, 0, -1) est la suite 3, 2, 1.
\end{description}

\begin{Quizz}
    Quelle suite donne l'appel de la fonction range suivante : range(1) ?
    \begin{itemize}
        \renewcommand{\labelitemi}{$\square$}
        \item 0
        \item 0, 1
        \item Rien
    \end{itemize}
\end{Quizz}
\begin{Quizz}
    Quelle suite donne l'appel de la fonction range suivante : range(-12, -8) ?
    \begin{itemize}
        \renewcommand{\labelitemi}{$\square$}
        \item -12, -11, -10, -9, -8
        \item -12, -11, -10, -9
        \item -11, -10, -9
    \end{itemize}
    Quelle suite donne l'appel de la fonction range suivante : range(20, 12, -4) ?
    \begin{itemize}
        \renewcommand{\labelitemi}{$\square$}
        \item 20, 16, 12
        \item 12, 16
        \item 20, 16
    \end{itemize}
\end{Quizz}


Maintenant que nous savons créer des itérables avec range, voyons comment combiner cela avec la boucle for.

\begin{center}{\textbf{\large Un cas simple pour s'échauffer :}}\end{center}
\begin{lstlisting}[style=pythonstyle, language=python]
for i in range(3) :
   print("Bonjour !")
\end{lstlisting}
Dans le code ci-dessus, la fonction range crée la suite 0, 1, 2. Il y a trois éléments dans la suite, on va donc exécuter le code de la boucle 3 fois. Le programme affiche alors 3 fois "Bonjour !". 
\begin{lstlisting}[style=pythonstyle, language=python]
Bonjour !
Bonjour !
Bonjour !
\end{lstlisting}

\begin{center}{\textbf{\large On monte le rythme :}}\end{center}
\begin{lstlisting}[style=pythonstyle, language=python]
for i in range(3, 5) :
   print(i)
\end{lstlisting}
Dans le code ci-dessus, la fonction range crée la suite 3, 4. Il y a deux éléments dans la suite, on va donc exécuter le code de la boucle 2 fois. Au premier tour de boucle, la variable i aura la première valeur de la suite donc 3 on affichera alors 3. Au deuxième tour de boucle, la variable i aura la deuxième valeur de la suite, donc 4 on affichera alors 4. Ce programme affiche donc :
\begin{lstlisting}[style=pythonstyle, language=python]
3
4
\end{lstlisting}

\begin{center}{\textbf{\large Un dernier un peu plus chaud :}}\end{center}
\begin{lstlisting}[style=pythonstyle, language=python]
somme=0
for i in range(4) :
   somme = somme + i
print(somme)
\end{lstlisting}
Dans le code ci-dessus, la fonction range crée la suite 0, 1, 2, 3. Il y a quatre éléments dans la suite, on va donc exécuter le code de la boucle 4 fois.
Ça commence à devenir compliqué, faisons un tableau d'exécution :
\begin{center}
\begin{tabular}{c|c c   | l}
     {\color{blue}Ligne} & $somme$ & $i$ & {\color{gray} Commentaires}\\ \hline
     {\color{blue}1} & 0  & $\times$ & {\color{gray}La variable somme est initialisé à 0, i n'existe pas.} \\
     {\color{blue}2} & | & 0 & {\color{gray} On parcourt le premier élément de la suite 0,1,2,3 donc i =0} \\
     {\color{blue}3} & 0+0=0 & | & {\color{gray}$somme = somme + i = 0 + 0 = 0$} \\
     {\color{blue}2} & | & 1 & {\color{gray}On parcourt le deuxième élément de la suite 0,1,2,3 donc i =1}\\
     {\color{blue}3} & 0 + 1 = 1 & | & {\color{gray}$somme = somme + i = 0 + 1 = 1$} \\
     {\color{blue}2} & | & 2 & {\color{gray}On parcourt le troisième élément de la suite 0,1,2,3 donc i =2}\\
     {\color{blue}3} & 1 + 2 = 3 & | & {\color{gray}$somme = somme + i = 1 + 2 = 3$}\\
     {\color{blue}2} & | & 3 & {\color{gray}On parcourt le quatrième élément de la suite 0,1,2,3 donc i =3}\\
     {\color{blue}3} & 3 + 3 = 6 & | & {\color{gray}$somme = somme + i = 3 + 3 = 6$}\\
     {\color{blue}2} & | & | & {\color{gray}On a parcouru tous les éléments, on sort de la boucle}\\
     {\color{blue}4} & \multicolumn{2}{r}{Le programme affiche 6} &\\
\end{tabular}
\end{center}
Ouf ! On y voit un peu plus clair ! À chaque tour de boucle, on a ajouté la valeur de $i$ à la valeur de $somme$, on a utilisé la variable $somme$ comme un~\motcle{accumulateur}.


\begin{Quizz}
    \begin{lstlisting}[style=pythonstyle, language=python]
somme_pair = 0
for i in range(0, 8, 2):
    somme_pair = somme_pair + i
print(somme_pair)
\end{lstlisting}
Compléter le tableau d'exécution du code ci-dessus :
\begin{center}
\begin{tabular}{c|c c   | l}
     {\color{blue}Ligne} & $somme\_pair$ & $i$ \\ \hline
     {\color{blue}1} & 0  & $\times$  \\
     {\color{blue}2} & | & 0   \\
     {\color{blue}3} & 0+0=0 & |  \\
     {\color{blue}2} & | &  & \\
     {\color{blue}3} & 0 + 2 = 2 & |  \\
     {\color{blue}2} & | & 4 \\
     {\color{blue}3} &  & | \\
     {\color{blue}2} &  & 6 \\
     {\color{blue}3} & 6 + 6 = 12 & | \\
     {\color{blue}2} &  &  \\
     {\color{blue}4} & \multicolumn{2}{r}{Le programme affiche } &\\
\end{tabular}
\end{center}
\end{Quizz}

\subsection{while}
La boucle while ("tant que" en anglais) permet de répéter un bloc de code  tant qu'une condition est vérifiée. L'avantage de la boucle while est donc qu'il ne faut pas connaitre à l'avance le nombre de fois où on devra répéter notre boucle.
\begin{lstlisting}[style=pythonstyle, language=python]
while condition :
    cette ligne sera répétée tant que la condition est vrai
    celle-ci aussi
celle-là ne sera executée qu une fois
\end{lstlisting}

Voyons un exemple :
\begin{lstlisting}[style=pythonstyle, language=python]
note_nsi = 0
while note_nsi != 20 :
    print("Je travaille et je m'amuse en codant !")
    note_nsi = note_nsi + 1
print("Je suis un As de la programmation !")
\end{lstlisting}
Dans le programme ci-dessus, on démarre avec $note\_nsi=0$ et on va travailler jusqu'à que la note soit égale à 20, à ce moment-là, on sort de la boucle while en étant un as de la programmation !
On fera donc 20 tours de boucle et on affichera 20 fois "Je travaille et je m'amuse en codant !".

\begin{Comprendre}
La boucle while ne s'arrête que quand la condition est fausse, ainsi, il est possible de faire des boucles infinies (des programmes qui ne s'arrêtent jamais).
On peut le faire exprès :
\begin{lstlisting}[style=pythonstyle, language=python]
temps = 0
while True :
    temps = temps + 1
\end{lstlisting}
Ou non !
\begin{lstlisting}[style=pythonstyle, language=python]
a = 15
while a != 20 :
    a = a + 2
\end{lstlisting}
\end{Comprendre}

\begin{Quizz}
\begin{lstlisting}[style=pythonstyle, language=python]
a = 14
while a % 6 != 0 :
    a = a - 1
\end{lstlisting}

Combien va-t-on faire de tour de boucle ?
\begin{itemize}
    \renewcommand{\labelitemi}{$\square$}
    \item 1
    \item 2
    \item 3
\end{itemize}
\end{Quizz}

% \newpage
